\documentclass[a4paper,10pt]{article}
\usepackage[francais]{babel}
\usepackage{url}

\selectlanguage{francais}
\input{letter-icube-UdS.sty}

\selectlanguage{francais}
\begin{document}


\name{St\'ephane \textsc{Genaud}}
\fonction{Professor}
\telephoneillkirch{+33 (0)3688 54542}
\email{genaud@unistra.fr}

\begin{letter}[Recommendation of Marc Fr\^incu for the CLOUD/NFV researcher
    position]
{to}
{Alcatel-Lucent Bell Labs}

Dear Madam/Sir, 

I first met Marc  Fr\^incu in 2007 when he was a student  on an internship in my
team at the  INRIA research center of Nancy,  France. I was at that  time on the
leave from my university on a researcher position. Although we have not directly
worked  together at  that time,  I  know from  my colleagues  that he  perfectly
achieved  his objectives  during  his  internship, which  led  my colleagues  to
propose him to apply for a \emph{CORDI}  PhD funding at INRIA. Marc preferred to
return  to Romania,  where  he  went on  with  a PhD  on  a  subject related  to
scheduling applications for Grids. This research  theme is close to my owns, and
I noticed  during the last  years the interesting  results published by  Marc in
conferences in the domain.

In  2012, I  learned  that Marc  was  searching  for a  PostDoc  position and  I
immediately proposed  him to  join us,  to work on  scheduling for  clouds. This
theme  perfectly  matched Marc's  professional  objectives.  Marc joined  us  on
October     2012,     to     participate     in     the     national     project
SONGS\footnote{http://infra-songs.gforge.inria.fr/}  I  am   involved  in.   The
objective  is to  enrich the  existing SimGrid  simulation toolkit  with an  API
dedicated to the  simulation of cloud-based systems. Simultaneously  to the work
on  simulation, we  carry  on theoretical  studies on  the  impact of  different
scheduling strategies  on different types  of workloads, that may  correspond to
certain  activity  types,  or  certain application  models  (e.g  submission  of
independent  jobs, or  workflows,  etc).  These  scheduling  algorithms will  be
candidate solutions  that could  benefit from  our cloud  simulator development.
Ideally, our  simulation results will  be compared to  those produced in  a real
software framework, such as the one under development in our team.

Marc has worked one year on this project (until september 2013). I must say Marc
has been  incredibly active and autonomous  during this period. We  have started
several articles,  always on Marc's  initiative, and all were  finally published
but one still under  review. During his stay, Marc has  shown an impressive body
of general  knowledge in the  field of grid and  cloud computing. We  have taken
this opportunity  to make a survey  of the existing literature  on brokering for
IaaS clouds.  Among other tasks, he has helped in the supervision of a PhD and a
Master student in the group, thereby showing a great teamwork spirit.


As a conclusion, I have been favourably impressed by the capabilities of Marc to
quickly identify and summarize relevant research issues, and then to write it in
a very  fluent style.  He  is both a  very autonomous, and  hard-working person.
Marc's stay (1 year) at or institution  has been very fruitful and allowed us to
achieve almost all  the numerous objective we assigned  ourselves.  The position
he  is applying  for at  the  Alcatel-Lucent Bell  labs center  to work  Network
Function Virtualization seems well in the scope of his expertise, given his work
on clouds and high availability.  Therefore,  I strongly recommend Marc for this
position.


\end{letter}
\begin{flushright}
\includegraphics[width=.30\textwidth]{signgenaud.jpg}
\end{flushright}
%\newpage
%\bibliographystyle{plain}
%\bibliography{julien_blaise_ED}


\end{document}
